% Track e6 -- The b-field on the temporal triple, c on the spatial triple, and the mass form T = e1e2e3.
% Verification: investigations/emergence/verify/e6_bfield_triples.py (claims e6-0 ... e6-7e).
% New references: sections/bib_e6.tex.
\section{\texorpdfstring{The $b$-field on the temporal triple, $c$ on the spatial triple}{The b-field on the temporal triple, c on the spatial triple}}
\label{e6:sec:main}

An earlier assignment places the arena of $c$ (of $c^2$) on the negative triple $\langle e_4,e_5,e_6\rangle$, of signature $(-,-,-)$, and the
arena of the $b$-field on the positive triple $\langle e_1,e_2,e_3\rangle$, of signature $(+,+,+)$, with the remark that the mass is unique in
$T=e_1e_2e_3$. In the conventions of \cite{ES17} the positive triple is the temporal triple $V^+$ and carries $\bB$, the negative triple is
the spatial triple $V^-$ and carries $\bE$ ($X=\sum_i(B_ie_i+E_ie_{i+3})$, Theorem~12.3 of \cite{ES17}), the clock is $\gamma^0=e_1$, the spatial
gamma matrices are $\gamma^a=e_{a+3}$, and $c$ is carried by the conversion $C\colon V^+\to V^-$, of which $C=c\,\id$ is the unique choice that
commutes with the diagonal rotations (Theorem~3.2 of \cite{ES17}; Proposition~2.7 of \cite{RE}). We therefore read the assignment as follows:
the $b$-field takes values in the temporal rotations $\so3_+=\operatorname{span}\{e_2e_3,e_3e_1,e_1e_2\}$, and $c^2$ is a symmetric part of $C$.
The trivector $T=\gamma^0J$ ($J=e_2e_3$) is the Dirac mass form, the unique antisymmetric kernel invariant under all fifteen bivectors
(Corollary~18.4 of \cite{ES17}), kept by $\mathrm{Spin}(3,3)$ at group level (Proposition~3.5 of \cite{RE}).

The answers, with their classes, are these. (1) As a form $T$ is invariant under all of $\so{3,3}$; as an operator it commutes exactly with
$\so3_+\oplus\so3_-$, and its commutant is $M_4(\C)$ (fixed; \cref{e6:sec:Tform}). (2) The temporal rotations keep the mass form and the rest
flow, and the kinetic term breaks them to $\uu1_J$, in proportion to the speed of the clock. They and the Pauli--G\"ursey $\su2$ of the massless
kinetic term are two diagonals of one $\so4$ meeting in $\R J$ (fixed; \cref{e6:sec:dirac}). (3) A temporal $b$-field is half a field-strength
rotation and half a counter-rotation. It keeps the rates, splits the lines into $\kappa$ and $s\pm\omega/2$, $s=(\kappa^2+\omega^2/4)^{1/2}$, cannot be
removed along a worldline, and tears the binding (fixed; \cref{e6:sec:binding}). (4) The binding root $G_C$ preserves
$\tfrac12(c^2|\bB|^2+|\bE|^2)$ iff $C=cR$, $R\in O(3)$. A temporal rotation keeps the symmetric factor $CC^{\!\top}$ on the spatial triple and moves
only the orthogonal factor (fixed; the reading ``$c^2$ on $(-,-,-)$, $b$ on $(+,+,+)$'' is interpretive; \cref{e6:sec:polar}). (5) Under a temporal
$b$-field the four-dimensional mass read by a fixed clock oscillates at fixed six-dimensional mass, unless the clock or the temporal energy is on
the axis (fixed; reading interpretive; \cref{e6:sec:clock}). (6) Between two constituents the temporal copy rotation rotates the graded masses at
fixed $M=(m_1^2+m_2^2)^{1/2}$, half of it creates quanta in pairs, and in the ungraded reading no transfer at fixed spectrum exists (fixed;
\cref{e6:sec:copies}). (7) At every rung the $\so{3n,3n}$-invariant forms are spanned by the temporal and the spatial volume; their symmetry has
period four in $n$, so the invariant Grassmann mass kernel is the temporal volume only for $n\equiv1\pmod4$ (fixed; \cref{e6:sec:ladder}).
\Cref{e6:sec:compare} compares the assignment with the $b$-field of Sections~5 and~9 of \cite{RE}. Identifiers e6-$n$ refer to the \texttt{PASS} lines of
\texttt{verify/e6\_bfield\_triples.py} (\cref{e6:tab:claims}).

\paragraph{Notation.} $S=\R^8$ with the realisation of equation~(1) of \cite{ES17}; $\omega_T=T=e_1e_2e_3$, $\omega_S=e_4e_5e_6=\gamma^1\gamma^2\gamma^3$,
$\Gamma=\omega_T\omega_S=e_1\cdots e_6$, $J=e_2e_3$, $\Omega=e_1e_4e_5e_6$ (Theorem~18.1 of \cite{ES17}); $K^\mu=-\gamma^0\gamma^\mu$ are the kinetic kernels
(Proposition~4.5 of \cite{ES17}) and $\mathfrak g_{\rm massless}=\operatorname{span}\{\Omega,J,e_2\Omega,e_3\Omega\}=\uu1_\Omega\oplus\su2_{\rm PG}$ is the
massless kinetic algebra (Proposition~5.5 of \cite{ES17}). The temporal generators are symmetric and the spatial ones antisymmetric
(Theorem~2.3 of \cite{ES17}). $\so3_\pm$ are the rotations of $V^\pm$, the six compact bivectors, and the nine bivectors $e_ie_{a+3}$ are the
boosts. A bivector $Y$ acts on vectors by $v\mapsto Y(v):=\tfrac12[Y,v]$. On $V=V^+\oplus V^-$ we use the blocks of \cite{RE}: $G=\bigl(\begin{smallmatrix}0&-\id\\ \id&0\end{smallmatrix}\bigr)$,
$G_C=\bigl(\begin{smallmatrix}0&-C^{-1}\\ C&0\end{smallmatrix}\bigr)$ (Proposition~2.7 of \cite{RE}), $N_C=\operatorname{diag}(C^{\!\top}C,\id)$ and
$E_C(X)=\tfrac12X^{\!\top}N_CX=\tfrac12(|C\bB|^2+|\bE|^2)$, with $E_c$ for $C=c\,\id$; $\lG\oplus\mathfrak p_G$ is the splitting of $\so{3,3}$ into the
parts commuting and anticommuting with $G$ (Proposition~5.3 of \cite{RE}). The interior shift is $S=b+\kappa G$ with phase-space generator $M$
(equation~(9) and Lemma~5.4 of \cite{RE}). Following Section~9 of \cite{RE}, the \emph{rates} are the spectrum of the $\eta$-self-adjoint part of
$S$ and the \emph{lines} are the frequencies of $S$ itself. On the tower $E_{(k,a)}$, $N$ and the ungraded generators $u_a=\id\otimes e_a$ are as in
\cref{e1:sec:main}; $C_b=-\tfrac12\sum_a\eta_{aa}E_{(1,a)}E_{(2,a)}$ is the copy rotation. Units are $\hbar=c=1$ unless $c$ is shown. The companion
note \emph{Two readings of $X$} \cite{TR} is cited by its own numbers.

\subsection{\texorpdfstring{$T$ as a form and as an operator}{T as a form and as an operator}}
\label{e6:sec:Tform}

\begin{proposition}[The two invariances of the mass form]
\label{e6:prop:Tform}
\begin{enumerate}[label=(\roman*),leftmargin=2.2em]
\item \emph{As a form.} For every bivector $Y$, $Y^{\!\top}T+TY=0$ and $Y^{\!\top}TY=-Y^2\,T$, so $(e^{sY})^{\!\top}Te^{sY}=T$ for every real $s$. The same holds for
$\omega_S$. The forms on $S$ invariant under $\so{3,3}$ are $\operatorname{span}\{T,\omega_S\}$, with $T$ antisymmetric and $\omega_S$ symmetric.
\item \emph{As an operator.} $T$ commutes with the six compact bivectors and anticommutes with the nine boosts. Hence $e^{sY}$ commutes with $T$ for
$Y\in\so3_+\oplus\so3_-$, while $e^{sY}Te^{-sY}=e^{2sY}T$ for a boost.
\item \emph{The commutant.} The commutant of $T$ in $M_8(\R)$ is spanned by the $32$ blades with an even number of spatial indices. It is
$\Cl(V^+)\cdot\Cl^0(V^-)$, a product of two commuting subalgebras with $\Cl(V^+)\cong\Cl(3,0)\cong M_2(\C)$, whose centre is $\operatorname{span}\{\id,T\}$, and
$\Cl^0(V^-)\cong\HH$. So it is $M_2(\C)\otimes_\R\HH\cong M_4(\C)$, of dimension $32$ and with centre $\operatorname{span}\{\id,T\}$, and it meets $\so{3,3}$ in
$\so3_+\oplus\so3_-\cong\so4$. The $32$ blades with an odd number of spatial indices anticommute with $T$.
\item $\Ad_T$ acts on $V$ as $\eta$: on the frame it is parity.
\end{enumerate}
\end{proposition}

\begin{proof}
By the blade rule (Lemma~2.1 of \cite{ES17}) a blade $e_I$ commutes with $T$ iff $3|I|-|I\cap\{1,2,3\}|$ is even, that is iff $|I\cap\{4,5,6\}|$ is even,
and anticommutes otherwise. Conjugation by $T$ is diagonal in the blade basis, so its commutant is spanned by the blades it contains, as in
\cref{e1:prop:copies}(ii). This gives (ii) and the first sentence of (iii). A temporal generator commutes with a spatial bivector, so
$\Cl(V^+)$ and $\Cl^0(V^-)$ commute, and their products are the $8\times4$ commuting blades. $T$ is central in $\Cl(V^+)$; $e_4e_5$, $e_5e_6$ and their product
$-e_4e_6$ square to $-\id$. The isomorphisms $\Cl(3,0)\cong M_2(\C)$ and $M_2(\C)\otimes_\R\HH\cong M_4(\C)$ are inherited; equivalently, $T^2=-\id$ makes
$S\cong\C^4$, and the commutant of a complex structure is $\mathfrak{gl}(4,\C)$. For (i): a compact $Y$ is antisymmetric with $Y^2=-\id$, so
$Y^{\!\top}T+TY=[T,Y]=0$ and $e^{sY}=\cos s+\sin s\,Y$; a boost is symmetric with $Y^2=\id$, so $Y^{\!\top}T+TY=\{T,Y\}=0$ and
$(\cosh s+\sinh s\,Y)^{\!\top}T(\cosh s+\sinh s\,Y)=(\cosh^2s-\sinh^2s)T$. $\omega_S$ has the same commutation pattern. The dimension two is
\cref{e6:thm:ladder}(i) at $n=1$. (iv) $T$ commutes with $e_1,e_2,e_3$ and anticommutes with $e_4,e_5,e_6$.
\end{proof}

Part (i) restates Corollary~18.4 of \cite{ES17} and Proposition~3.5(i) of \cite{RE} (fixed there), with the symmetric partner $\omega_S$ of
Section~28 of \cite{ES17}. Part (iv) restates Proposition~2.11(i) of \cite{RE}. Parts (ii) and (iii) are fixed (e6-1a to e6-1c); that $\omega_T$ and
$\omega_S$ are fixed by all six compact rotations is noted after Theorem~18.1 of \cite{ES17}. The two invariances are one statement read twice.
Because the compact bivectors are antisymmetric and the boosts symmetric, invariance of the form is commutation with the first and
anticommutation with the second. As a form, $T$ breaks nothing in $\so{3,3}$. As an operator it singles out the maximal compact subalgebra, the
rotations of the two triples, and among these it does not distinguish the temporal from the spatial ones. Two things distinguish $T$ from
$\omega_S$: as a Grassmann kernel $T$ is antisymmetric and $\omega_S$ symmetric, and as an operator $T^2=-\id$ is a complex structure while
$\omega_S^2=\id$ gives a projector (Theorem~18.1 of \cite{ES17}).

\subsection{The temporal rotations on the Dirac data}
\label{e6:sec:dirac}

The Dirac data of one constituent are the kinetic kernels $K^\mu$, the rest Hamiltonian $H(0)=m\gamma^0=me_1$ of the real first-order form,
and the rest flow $e^{-m\lambda T}$, with $JH(0)=mT$ (Section~4.4 of \cite{ES17}; Propositions~3.7 and~4.9 of \cite{RE}).

\begin{proposition}[What the temporal rotations do to the Dirac data]
\label{e6:prop:dirac}
Let $Y=\alpha J+\beta\,e_3e_1+\gamma\,e_1e_2\in\so3_+$.
\begin{enumerate}[label=(\roman*),leftmargin=2.2em]
\item \emph{Kinetic forms.} $Y^{\!\top}K^0+K^0Y=0$ and $Y^{\!\top}K^a+K^aY=2\,Y(\gamma^0)\,\gamma^a$, where $Y(\gamma^0)=\beta e_3-\gamma e_2$ is the velocity
of the clock. Hence $\sum_a\|Y^{\!\top}K^a+K^aY\|^2=96(\beta^2+\gamma^2)$ (Frobenius norm), and $Y$ maps the kinetic forms into their span,
$Y^{\!\top}K^\mu+K^\mu Y=\Lambda^\mu{}_\nu K^\nu$, iff $\beta=\gamma=0$: $\so3_+\cap(\lint\oplus\R J)=\R J$.
\item \emph{Rest Hamiltonian.} $e^{\varphi e_1e_2}e_1e^{-\varphi e_1e_2}=\cos2\varphi\,e_1-\sin2\varphi\,e_2$, and likewise with $e_1e_3$ and $e_3$. The stabiliser of
$H(0)$ in $\so3_+$ is $\R J$. A rotor of $\mathrm{Spin}(V^+)$ that carries $e_1$ to a clock $n$ carries $J$ to $J(n)=Tn$, $K^a$ to $-n\gamma^a$ and
$H(0)$ to $mn$, and it fixes $K^0=-\id$ and $T$. Moreover $nJ(n)=J(n)n=T$.
\item \emph{Rest flow.} $e^{-m\lambda T}$ commutes with $\mathrm{Spin}(3)_+\times\mathrm{Spin}(3)_-$. The spin lift $\tfrac12e_je_k$ of $\so3_+$ has Casimir
$-\tfrac34\id$ on $S$, so $S$ carries isospin $\tfrac12$, and $e^{\pi e_je_k}=-\id$. The clock rotations $e_1e_2$ and $e_3e_1$ anticommute with $J$ and commute
with $T$. Their half turn $\Ad_{e_1e_2}=\Ad_{e_1}\Ad_{e_2}$ maps $(e_1,J,T,\gamma^a)$ to $(-e_1,-J,T,\gamma^a)$.
\item \emph{Field octants.} With the field-octant map $\Phi\colon V\to S$ of Theorem~4.7 of \cite{TR}, $-T\Phi=\Phi G$, so
$e^{-m\lambda T}\Phi=\Phi\,e^{m\lambda G}$: on the six field octants the Dirac rest flow is the duality rotation at angular frequency $m$.
\end{enumerate}
\end{proposition}

\begin{proof}
(i) $Y$ is antisymmetric and commutes with each $\gamma^a$ (a temporal bivector and a spatial generator), so
$Y^{\!\top}K^a+K^aY=[K^a,Y]=-[e_1,Y]\gamma^a=2Y(e_1)\gamma^a$, and $K^0=-\id$. With $v=Y(e_1)$ temporal, $(v\gamma^a)^{\!\top}(v\gamma^a)=|v|^2\id$, so each term has
squared norm $4\cdot8|v|^2$. The covariance statement intersects the seven-dimensional solution space $\lint\oplus\R J$ of Proposition~3.5(ii) of
\cite{RE} with $\so3_+$; it is also solved directly (e6-2a). (ii) $e_1$ anticommutes with $e_1e_2$, so $e^{\varphi e_1e_2}e_1e^{-\varphi e_1e_2}=e^{2\varphi e_1e_2}e_1$. The
stabiliser condition is $Y(e_1)=0$. A rotor $R$ of $\mathrm{Spin}(V^+)$ commutes with $T$ and with every $\gamma^a$, and $J=Te_1$, so $\Ad_RJ=Tn$ and
$\Ad_RK^a=-n\gamma^a$. Since $n$ commutes with $T$ and $n^2=\id$, $nJ(n)=J(n)n=T$. (iii) $T$ commutes with the compact bivectors by
\cref{e6:prop:Tform}(ii); the Casimir is $\tfrac14\sum Y_j^2$ with $Y_j^2=-\id$; the rest is the blade rule. (iv) Iterate $-T\Phi=\Phi G$.
\end{proof}

Part (i) is fixed (e6-2a); its covariance half restates Proposition~3.5(ii) of \cite{RE}. Part (ii) is fixed (e6-2b, exact at $\cos\varphi=\tfrac35$
and at the clock $n=-\tfrac7{25}e_1-\tfrac{24}{25}e_2$, where (i) holds with $Y(n)$ in place of $Y(\gamma^0)$). Its transport of the data restates
Proposition~3.4 and Theorem~3.5 of \cite{TR}, by which the Dirac mass form $\gamma^0J(n)=T$ is the same for every clock while $J(n)$ turns with $n$.
Part (iii) is fixed (e6-2c). Part (iv) restates Proposition~4.12(i) of \cite{TR}, re-verified with $\Phi$ built from the Fock vacuum (e6-2c).

So the temporal $SU(2)$ is a symmetry of the mass form and of the rest flow that the kinetic term breaks to $\uu1_J$. That reverses the usual
order, kinetic term symmetric and mass breaking, as Section~3.3 of \cite{RE} found for $\mathrm{Spin}(3,3)$. The broken directions are the two clock
rotations: the kinetic term pairs one temporal direction with the spatial triple, and the breaking is the speed $|Y(n)|$ of the clock. They flip the
charge, and their half turn acts as $\mathsf{CP}$ ($\Ad_{e_2}$ is the charge conjugation of Theorem~5.1 of \cite{ES17}, $\Ad_{e_1}$ the parity), reversing
the sense of the rest rotation as in Proposition~3.8(iii) of \cite{RE}. Interpretively, the temporal $SU(2)$ joins particle and antiparticle through
the sphere of clocks at fixed mass form.

\begin{proposition}[The mass symmetry and the kinetic symmetry]
\label{e6:prop:twosu2}
Let $P_\pm=\tfrac12(\id\pm\omega_S)$.
\begin{enumerate}[label=(\roman*),leftmargin=2.2em]
\item $\so3_+\cap\mathfrak g_{\rm massless}=\R J$.
\item $e_2\Omega=-(e_1e_2)\,\omega_S$ and $e_3\Omega=(e_3e_1)\,\omega_S$. $\omega_S$ commutes with $\so3_+$, and $P_\pm$ are orthogonal projectors of rank four.
\item $\so3_+=\{YP_++YP_-\}$ and $\su2_{\rm PG}=\{YP_++(JYJ^{-1})P_-\}$, $Y\in\so3_+$: both are graphs of automorphisms of $\so3$ over the two
eigenspaces of $\omega_S$, the identity and the half turn $\Ad_J$ about $J$. The Lie algebra they generate is $\so3_+P_+\oplus\so3_+P_-\cong\so4$, and they
meet in the fixed algebra of $\Ad_J$, which is $\R J$.
\item $g=P_++JP_-=\exp(\tfrac\pi2JP_-)$ is orthogonal and $g\,\so3_+\,g^{-1}=\su2_{\rm PG}$; $2g=\id+J+\omega_S-J\omega_S$ has components of grades $0,2,3,5$,
so $g\notin\mathrm{Pin}(3,3)$.
\item $\so3_+$ preserves $T$ and the kinetic term breaks it to $\R J$; $\su2_{\rm PG}$ preserves the kinetic forms and the mass breaks it to $\R J$.
\end{enumerate}
\end{proposition}

\begin{proof}
(ii) $e_2\Omega=e_2e_1e_4e_5e_6=-e_1e_2\,\omega_S$, and similarly for $e_3$. $\omega_S$ is disjoint from each temporal bivector, and the grades $3$ and $2$
give a commuting sign. (iii) Substitute $\omega_S=P_+-P_-$. $J$ commutes with $J$ and anticommutes with $e_1e_2$ and $e_3e_1$, so
$\Ad_J$ is the half turn about $J$. A bracket of two graphs of automorphisms $\phi_1,\phi_2$ contains $[Y,Y']P_++[\phi_1Y,\phi_2Y']P_-$, and the
generated algebra is computed (e6-2d). (iv) $(JP_-)^2=-P_-$ and $gYg^{-1}=YP_++JYJ^{-1}P_-$, because $P_\pm$ commute with $Y$ and $J$. (i) follows from (iii),
and (v) is \cref{e6:prop:dirac}(i) together with Proposition~3.6(i)--(ii) of \cite{RE}.
\end{proof}

This is fixed (e6-2d); (v) restates Proposition~5.5 of \cite{ES17} and Proposition~3.6 of \cite{RE}. The comparison has a sharp answer.
The two $\su2$'s are not the same, and neither contains the other. They are the untwisted and the $\Ad_J$-twisted diagonals of $\so3_+P_+\oplus\so3_+P_-$,
and they share the Cartan direction $J$. One is the symmetry of the mass that the kinetic term breaks, the other the symmetry of the kinetic term
that the mass breaks. Both leave the same $\uu1_J$ \cite{Pauli57,Gursey58}. The orthogonal map that exchanges them is not in the Clifford group.

\subsection{A temporal \texorpdfstring{$b$}{b}-field and the binding root}
\label{e6:sec:binding}

\begin{proposition}[A $b$-field on the temporal triple]
\label{e6:prop:binding}
Let $b=(R,0)$ on $V^+\oplus V^-$, with $R=[\mathbf r]_\times\neq0$, $\omega=|\mathbf r|$, $\kappa>0$ and $S=b+\kappa G$.
\begin{enumerate}[label=(\roman*),leftmargin=2.2em]
\item $b=b_l+b_p$ with $b_l=\tfrac12(R,R)\in\lG$, a field-strength rotation, and $b_p=\tfrac12(R,-R)\in\mathfrak p_G$, a counter-rotation; they commute. In
the notation of equation~(11) of \cite{RE}, $\mathbf a=\mathbf c=\mathbf r/2$, so $\mathbf a\cdot\mathbf c=\omega^2/4$.
\item $e^{\theta b}Ge^{-\theta b}=G_{C(\theta)}$ with $C(\theta)=e^{-\theta R}$. $G_C^2=-\id$ for every $C$, and $G_C$ is $\eta$-self-adjoint iff $C$ is
orthogonal; so $\kappa G_{C(\theta)}$ stays in the self-adjoint summand and the rates, $\pm i\kappa$, are unchanged. $[b,G]=2b_pG\neq0$, so $b$ is not removable along a worldline.
\item $\det(\lambda-S)=(\lambda^2+\kappa^2)\bigl(\lambda^4+(2\kappa^2+\omega^2)\lambda^2+\kappa^4\bigr)$. The lines are $\kappa$, on the plane
$\langle\hat{\mathbf r},G\hat{\mathbf r}\rangle$, and $s\pm\omega/2$, with $s=(\kappa^2+\omega^2/4)^{1/2}$ and $(s+\omega/2)(s-\omega/2)=\kappa^2$; at $(\kappa,\omega)=(2,3)$
they are $1,2,4$.
\item Not every solution is bounded. At $(\kappa,\omega)=(2,3)$ and $(6,5)$, $M$ has Jordan blocks exactly at the frequencies $s\pm\omega/2$ and is
semisimple at $\kappa$.
\item On spinors $e_je_k=\tfrac12(e_je_k-e_{j+3}e_{k+3})+\tfrac12(e_je_k+e_{j+3}e_{k+3})$ for temporal $j,k$. The first summand, a paired rotation, commutes with
$N$; the second changes $N$ by $\pm2$ only.
\end{enumerate}
\end{proposition}

\begin{proof}
(i) $G$ commutes with $\operatorname{diag}(R,R)$ and anticommutes with $\operatorname{diag}(R,-R)$ in the $(\bB,\bE)$ blocks. (ii)
$e^{\theta b}=\operatorname{diag}(e^{\theta R},\id)$, and $\operatorname{diag}(Q,\id)\,G\,\operatorname{diag}(Q^{-1},\id)=G_{Q^{-1}}$. $G_C^\dagger=\eta G_C^{\!\top}\eta$ has blocks $-C^{\!\top}$ and
$C^{-\top}$, which equal those of $G_C$ iff $C^{\!\top}=C^{-1}$. The rates are the spectrum of the self-adjoint part, a tensor (Theorem~9.3(iii) of
\cite{RE}). Removability fails by Proposition~5.9 of \cite{RE}. (iii) $b_l$ commutes with $b_p$ and with $G$. By Larmor's theorem (Proposition~5.7 of
\cite{RE}) the substitution $X=e^{-b_l\tau}Y$ removes $b_l$ and shifts the exponents by the eigenvalues $0,\pm i\omega/2$ of $-b_l$. On the support of $b_p$ the
remaining shift $b_p+\kappa G$ has the frequency $(\kappa^2+\omega^2/4)^{1/2}$ of Theorem~5.8(ii) of \cite{RE}, and on its kernel $\langle\hat{\mathbf r},G\hat{\mathbf r}\rangle$ the
frequency $\kappa$. The characteristic polynomial is also computed symbolically in $\mathbf r$. (iv) By Theorem~5.10(i) of \cite{RE}, boundedness needs
$\mathbf a\cdot\mathbf c=0$. The Jordan structure is that of $b_p+\kappa G$ (Theorem~5.8(ii) of \cite{RE}), transported by the bounded frame $e^{-b_l\tau}$, and is
computed exactly. (v) With $e=a^++a^-$ and $e'=a^+-a^-$, the sum is $2(a^+_ja^+_k+a^-_ja^-_k)$ and the difference $2(a^+_ja^-_k+a^-_ja^+_k)$.
\end{proof}

This is fixed (e6-3a to e6-3e); the statement on the rates restates Theorem~9.3(iii) of \cite{RE}. A $b$-field on the temporal triple alone is
therefore half binding-compatible and half binding-odd. Its paired half is the compact part of $\lG$, a Zeeman term removable by a rotating frame
\cite{Larmor1897}. Its counter half lies in $\mathfrak p_G$. The latter is the gauge-covariant remainder $D_\tau G=[b,G]$, it moves the binding root
to the root $G_{e^{-\theta R}}$ of a rotated conversion, and it tears the binding, with the screened frequency $s$ and Jordan blocks. The rate
$\kappa=c/\ell$ is never touched. In \cref{e6:prop:binding}(iii) the two halves act independently, because $\mathbf a\parallel\mathbf c$: the lines are
the screened frequency $\sqrt{\kappa^2+|\mathbf c|^2}$ shifted by the Zeeman amounts $\pm|\mathbf a|$. Theorem~5.10 of \cite{RE} needed
$\mathbf a\perp\mathbf c$ for screening; a field on one triple has the opposite geometry. On spinors the same split is the split of
\cref{e1:prop:Nbiv}(ii): the paired half hops quanta within a constituent and the counter half creates or annihilates them in pairs, so the
temporal rotations, $J$ included, do not conserve $N$ (\cref{e1:prop:Nlorentz}(ii); \cref{e2:lem:N}).

\subsection{\texorpdfstring{$c$}{c} on the spatial triple: the polar form of the conversion}
\label{e6:sec:polar}

\begin{proposition}[Conversions and the temporal rotations]
\label{e6:prop:polar}
Let $C\colon V^+\to V^-$ be invertible, $z_C=C\bB+i\bE$ and $q_C=z_C\cdot z_C$.
\begin{enumerate}[label=(\roman*),leftmargin=2.2em]
\item $G_C^2=-\id$, $G_C$ acts on $z_C$ as multiplication by $i$, and $e^{\theta G_C}$ preserves $E_C$, which depends on $C$ only through $C^{\!\top}C$.
\item $G_C^{\!\top}N_c+N_cG_C=\bigl(\begin{smallmatrix}0&C^{\!\top}-c^2C^{-1}\\ C-c^2C^{-\top}&0\end{smallmatrix}\bigr)$, so $G_C$ preserves
$E_c=\tfrac12(c^2|\bB|^2+|\bE|^2)$ iff $C^{\!\top}C=c^2\id$, that is iff $C=cR$ with $R\in O(3)$.
\item For $C=cR$, the elements of $\so3\oplus\so3$ commuting with $G_C$ are the twisted diagonal rotations $(Q,RQR^{-1})$. $G_{cR}=\Ad_{(R^{-1},\id)}G_c=\Ad_{(\id,R)}G_c$
for $R\in SO(3)$, and $\Ad_{(R_+,R_-)}G_c=G_{cR_-R_+^{-1}}$ for $(R_+,R_-)\in O(3)\times O(3)$. The Lie algebra preserving $\eta$ and $E_c$ is
$\so3\oplus\so3$, and adding $[Y,G_c]=0$ leaves the diagonal $\so3$. Hence $\{G_C:C^{\!\top}C=c^2\id\}$ is one orbit of $O(3)\times O(3)$, of dimension three.
\item A temporal rotation $(R_+,\id)$ carries $G_C$ to $G_{CR_+^{-1}}$ and $N_C$ to $N_{CR_+^{-1}}$. It conjugates $C^{\!\top}C$ to $R_+C^{\!\top}CR_+^{-1}$ and fixes
$CC^{\!\top}$; in the polar form $C=(CC^{\!\top})^{1/2}U$ it moves only $U$. A spatial rotation $(\id,R_-)$ carries $C$ to $R_-C$, fixes $C^{\!\top}C$ and
moves only $U'$ in $C=U'(C^{\!\top}C)^{1/2}$.
\item For $C=cR$: $E_c$ and $\Re q_C=c^2|\bB|^2-|\bE|^2$ do not depend on $R$, while $\Im q_C=2c\,(R\bB)\cdot\bE$ does. At $c=1$, $\bB=e_1$, $\bE=e_2$ and $R$ the
rotation through $\theta$ about $e_3$, $q_C=2i\sin\theta$: the same field point is null for $\theta=0$ and at rest for $\theta=\tfrac\pi2$, where
$E_0=E_6=1$.
\end{enumerate}
\end{proposition}

\begin{proof}
(i) $G_C^2$ has blocks $-C^{-1}C$ and $-CC^{-1}$. On $z_C$, $G_CX$ has $C\bB'=-\bE$ and $\bE'=C\bB$, so $z_C(G_CX)=iz_C(X)$, and $|z_C|^2=2E_C$. (ii) is a
block computation. Its vanishing means $C^{\!\top}C=c^2\id$, and then $C/c$ is orthogonal. (iii) The lower-left block of $[G_C,\operatorname{diag}(Q,Q')]$ is
$CQ-Q'C$. The conjugation formulas are block products. An $\eta$-isometry preserving the positive form $N_c$ preserves the eigenspaces $V^\pm$ of
$\eta^{-1}N_c=\operatorname{diag}(c^2,-1)$, and is orthogonal on each. The stabiliser dimensions are exact null spaces (e6-4c). (iv)
$(CR_+^{-1})(CR_+^{-1})^{\!\top}=CC^{\!\top}$. (v) Expand $z_C\cdot z_C$; in the example $R\bB=(\cos\theta,\sin\theta,0)$.
\end{proof}

This is fixed (e6-4a to e6-4e). Parts (i)--(iii) extend Proposition~2.7 of \cite{RE} and Theorem~3.2 of \cite{ES17} from $C=c\,\id$ to the orbit. They are
elementary linear algebra: Schur's lemma applied to a twisted diagonal action gives $C=cR$. Part (v) is the passive form of Proposition~2.14(ii) of
\cite{RE}, where the anti-diagonal part of a compact transport turns a null field into one at rest at fixed $E_6$.

\begin{remark}[What ``$c^2$ on $(-,-,-)$, $b$ on $(+,+,+)$'' fixes]
\label{e6:rem:csym}
\begin{enumerate}[label=(\alph*),leftmargin=2.2em]
\item \emph{Fixed.} The complete invariant of a conversion under the temporal rotations is $CC^{\!\top}$, a symmetric form on the spatial triple:
$C$ and $C'$ lie in one orbit of the orthogonal group of the temporal triple iff $C'C'^{\!\top}=CC^{\!\top}$ (\cref{e6:prop:polar}(iv)). A $b$-field on the temporal triple can change
anything about the conversion except this. The paper's $c$ is the isotropic case $CC^{\!\top}=c^2\id$, which is also $C^{\!\top}C=c^2\id$, and that
is exactly the case in which the binding root preserves the weighted energy (\cref{e6:prop:polar}(ii)). In this precise sense $c^2$ sits on the
spatial triple and the temporal $b$-field cannot reach it. What the temporal $b$-field moves is the orthogonal factor $R$. The energy $E_c$ does not see
$R$; the binding root $G_{cR}$ (which spatial direction is paired with which temporal one), the pseudoscalar $\Im q_C$ and the rest energy
$\tfrac12|q_C|$ do.
\item \emph{Fixed: a mirror.} With the triples exchanged the statement still holds: a spatial $b$-field keeps $C^{\!\top}C$ on the temporal triple and
moves the other orthogonal factor. The conversion alone therefore does not decide which triple carries $c^2$ and which the $b$-field. The Dirac data
decide it, through the clock in $V^+$ and the kinetic term (\cref{e6:prop:dirac}): in the interior reading $\so3_-\subset\lint$ is part of the
Lorentz algebra, while $\so3_+$ meets $\lint\oplus\R J$ only in $J$. The frame reading does not decide it, since its Lorentz algebra $\lG$ contains
only the diagonal rotations.
\item \emph{Relation to the clause ``$c^2$ is the symmetric part of $E=mc^2$''.} Section~2.5 of \cite{RE} gives the clause its exact forms: $c$ is
the unique isotropic conversion, its weight in $E_6$ a unit convention (Propositions~2.7 and~2.8 there), and the shift splits as
$\End(V)=\so{3,3}\oplus\{Y^\dagger=Y\}$ with $\kappa G$ in the self-adjoint summand, $(\kappa G)^2=-c^2/\ell^2$, and the connections in $\so{3,3}$ (Section~5
there). The polar decomposition of the conversion is a further, compatible form: its symmetric factor carries $c^2$ and is invariant under the
temporal $b$-field, its orthogonal factor is moved by it, and for every orthogonal $C$, $\kappa G_C$ stays in the self-adjoint summand with the same
spectrum (\cref{e6:prop:binding}(ii)). It adds nothing else: it does not fix $c$, does not make $c$ a symmetry (the stabiliser $O(3)\times O(3)$ of $\eta$
and $E_c$ is not a Lorentz group), and selects no rest energy. A constant $R$ is a choice of identification of the triples (Theorem~3.2 of
\cite{ES17} selects $R=\id$ once a common $SO(3)$ is fixed); only a $\tau$-dependent $R(\tau)$, a temporal $b$-field, has gauge-invariant content, its
$D_\tau G\neq0$. Reading the symmetric factor as ``$c^2$'' and the orthogonal factor as the place of the $b$-field is interpretive.
\end{enumerate}
\end{remark}

\subsection{The six-momentum read by a clock}
\label{e6:sec:clock}

Let $P=\sum_i(\mathsf E_i/c)e_i+\sum_ap_ae_{a+3}$ with $\eta(P,P)=\mu^2c^2$, so $\mu^2c^4=|\boldsymbol{\mathsf E}|^2-|\bp|^2c^2$, as in Proposition~2.10 of \cite{RE}.
For a clock $n$ the energy read is $E_n=n\cdot\boldsymbol{\mathsf E}$ and $m_4(n)^2c^4=E_n^2-|\bp|^2c^2=\mu^2c^4-|n\times\boldsymbol{\mathsf E}|^2$ (Proposition~3.13 of
\cite{TR}). A temporal $b$-field of constant rate $\Omega$ about a unit axis $\hat{\mathbf e}$ acts on $P$ by $\boldsymbol{\mathsf E}(t)=R_{\hat{\mathbf e}}(\Omega t)\boldsymbol{\mathsf E}$,
with $\bp$ fixed.

\begin{proposition}[The four-dimensional mass under a temporal $b$-field]
\label{e6:prop:clock}
\begin{enumerate}[label=(\roman*),leftmargin=2.2em]
\item The flow keeps $\mu$ and $|\boldsymbol{\mathsf E}|$ constant.
\item For a fixed clock $n$, $E_n(t)=A+B\cos(\Omega t-\phi)$, where $A=(n\cdot\hat{\mathbf e})(\hat{\mathbf e}\cdot\boldsymbol{\mathsf E})$ and
$B=|n_\perp|\,|\boldsymbol{\mathsf E}_\perp|$ with $n_\perp$ and $\boldsymbol{\mathsf E}_\perp$ the parts orthogonal to $\hat{\mathbf e}$. Hence $m_4(n)^2c^4=\mu^2c^4-|\boldsymbol{\mathsf E}|^2+E_n(t)^2$ oscillates. Its maximum is
$\mu^2c^4-|\boldsymbol{\mathsf E}|^2+(|A|+B)^2\le\mu^2c^4$. Its minimum is $\mu^2c^4-|\boldsymbol{\mathsf E}|^2+(|A|-B)^2$ if $|A|\ge B$ and
$\mu^2c^4-|\boldsymbol{\mathsf E}|^2$ otherwise. It is constant iff $B=0$, that is iff $n\parallel\hat{\mathbf e}$ or $\boldsymbol{\mathsf E}\parallel\hat{\mathbf e}$.
\item The rotation about the clock, whose spin lift is proportional to $J(n)$ (Proposition~3.4(iii) of \cite{TR}), keeps $E_n$. On the reduced shell
$\boldsymbol{\mathsf E}\parallel n$ of Theorem~30.1 of \cite{ES17} (Proposition~3.10 of \cite{TR}) the flow preserves the constraint iff $\hat{\mathbf e}\parallel n$.
\item With $c=1$, $\boldsymbol{\mathsf E}=(2,3,6)$ and $\bp=(0,0,3)$, so that $\mu^2=40$: the rotation about $e_3$ read by $e_1$ gives $E_n=2\cos\Omega t-3\sin\Omega t$ and
$m_4^2\in[-9,4]$; read by $\widehat W$ it gives $m_4^2\in[\tfrac{35}3-4\sqrt{26},\tfrac{35}3+4\sqrt{26}]$. The rotation about $\widehat W$ read by $\widehat W$ keeps
$m_4^2=\tfrac{94}3$ while the relative-time energies $\boldsymbol{\mathsf E}_\perp$ turn.
\end{enumerate}
\end{proposition}

\begin{proof}
Rodrigues' formula gives $n\cdot\boldsymbol{\mathsf E}(t)=A+(n\cdot\boldsymbol{\mathsf E}-A)\cos\Omega t+n\cdot(\hat{\mathbf e}\times\boldsymbol{\mathsf E})\sin\Omega t$, and the squared
amplitude $(n\cdot\boldsymbol{\mathsf E}-A)^2+(n\cdot(\hat{\mathbf e}\times\boldsymbol{\mathsf E}))^2$ equals $(|n|^2-(n\cdot\hat{\mathbf e})^2)(|\boldsymbol{\mathsf E}|^2-(\hat{\mathbf e}\cdot\boldsymbol{\mathsf E})^2)$
(Lagrange's identity in the plane orthogonal to $\hat{\mathbf e}$). $|A|+B\le|\boldsymbol{\mathsf E}|$ by the Cauchy--Schwarz inequality. The rest is
substitution; the reduced-shell statement says that a vector stays on the line $\R n$ under all rotations about $\hat{\mathbf e}$ iff the line is the axis.
\end{proof}

This is fixed (e6-5a, symbolic for three rational axes; e6-5b, exact). The reading needs care. ``The four-dimensional mass oscillates'' holds only
if the clock is held fixed while the temporal $b$-field turns the temporal momentum. Equivalently, the clock turns against a fixed momentum along a
great circle of the clock sphere of Proposition~3.13 of \cite{TR}. Read in the sense of Remark~3.14 of \cite{TR}, the energy read by the clock is
the projection of the undivided temporal energy, and the $b$-field moves energy periodically between the clock and the relative times. This is the
dynamical form of the obstruction of Proposition~2.10 of \cite{RE}: splitting the temporal energy lowers the mass. The six-dimensional mass is
constant throughout. The rotation about the clock is the charge $J(n)$ and changes nothing; the two clock rotations are what make $m_4$ oscillate, and
they are also what the kinetic term breaks (\cref{e6:prop:dirac}). On the reduced shell, the gauge-fixed reading of \cite{ES17}, a temporal $b$-field
transverse to the clock is not admissible at fixed clock.

\subsection{Between constituents}
\label{e6:sec:copies}

Write $C^{12}_+=-\tfrac12\sum_{i\le3}E_{(1,i)}E_{(2,i)}$ and $C^{12}_-=\tfrac12\sum_{a\ge4}E_{(1,a)}E_{(2,a)}$ for the temporal and the spatial parts of the copy rotation,
$C_b=C^{12}_++C^{12}_-$, and $C_\times=C^{12}_+-C^{12}_-$. We drop the superscript.

\begin{proposition}[Temporal and spatial copy rotations]
\label{e6:prop:copyrot}
\begin{enumerate}[label=(\roman*),leftmargin=2.2em]
\item $C_\pm$ are antisymmetric and commute, and each is a sum of three commuting terms of square $-\tfrac14\id$. Hence $e^{2\pi C_\pm}=-\id$ and $e^{2\pi C_b}=\id$.
The eigenvalues of $C_\pm$ are $\pm\tfrac i2$ ($24$ times each) and $\pm\tfrac{3i}2$ ($8$ times each).
\item On vectors neither $C_+$ nor $C_-$ commutes with $G_2$, $C_b$ does, and $C_\times$ is compact and anticommutes with $G_2$.
\item $C_+=\tfrac12C_b+\tfrac12C_\times$. $C_b$ conserves $N$ and $C_\times=-\sum_i(a^+_{1,i}a^+_{2,i}+a^-_{1,i}a^-_{2,i})$ changes it by $\pm2$, one quantum in each
constituent. From the Fock vacuum, $e^{\theta C_+}|0\rangle$ has weight $\binom3k\cos^{2(3-k)}\tfrac\theta2\,\sin^{2k}\tfrac\theta2$ in $N=2k$, none in odd $N$,
and $N^{(1)}=N^{(2)}$ throughout. So $\langle N\rangle=3(1-\cos\theta)$, and $\theta=\pi$ fills all six modes.
\item As a constant field between two binding frames, $S=\kappa G_2+\beta V_+$, where $V_+$ is the vector form of $C_+$, has
$\det(\lambda-S)=\bigl(\lambda^4+(2\kappa^2+\beta^2)\lambda^2+\kappa^4\bigr)^3$. The rates are $\kappa$ and the lines are $s\pm\beta/2$. At $(\kappa,\beta)=(2,3)$, $M$ has Jordan blocks
at both lines.
\end{enumerate}
\end{proposition}

\begin{proof}
(i) $E_{(1,i)}E_{(2,i)}$ squares to $-\id$ for every $i$, and the terms for distinct $i$ involve four distinct anticommuting generators. Then
$e^{2\pi C_\pm}=\prod(\cos\pi\mp\sin\pi\,E_{(1,i)}E_{(2,i)})=(-1)^3$. The multiplicities follow from the minimal polynomial $(\lambda^2+1)(\lambda^2+9)$ of
$2C_\pm$ and $\Tr(2C_\pm)^2=-192$. (ii) $V_\pm$ maps $e^{(1)}_a$ to $e^{(2)}_a$ on the temporal (respectively spatial) directions, and $G_2$ exchanges the two
kinds; $V_+-V_-=\varepsilon\otimes\eta$. (iii) As in \cref{e6:prop:binding}(v): $E_{(1,i)}E_{(2,i)}\pm E_{(1,i+3)}E_{(2,i+3)}$ is a pair term or a hopping term. Since
$[C_b,C_\times]=0$ and $C_b|0\rangle=0$, $e^{\theta C_+}|0\rangle=e^{\theta C_\times/2}|0\rangle$, a product of three rotations in the planes
$\langle|0\rangle,a^+_{1,i}a^+_{2,i}|0\rangle\rangle$; this is checked exactly at $\cos\tfrac\theta2=\tfrac35$. (iv) As in \cref{e6:prop:binding}(iii), with
$\varepsilon\otimes\id_3$ in place of $[\hat{\mathbf r}]_\times$.
\end{proof}

Part (i) restates Proposition~9.6(i) of \cite{RE}; the rest is fixed (e6-6a, e6-6c, e6-6e). Part (iii) is the compact counterpart of the copy boost
of \cref{e2:prop:copyboost}. The pairs emerge unitarily, with probability $\binom3k\cos^{2(3-k)}\tfrac\theta2\sin^{2k}\tfrac\theta2$ in place of the hyperbolic
weights there. As \cref{e2:sec:quanta} says, compactness does not protect the number of quanta; commutation with the binding root does, and $C_+$
fails it. Part (iv) puts $C_+$ outside the binding-compatible class of Theorem~9.3 of \cite{RE}: as a constant field it tears the binding of two
manifolds as a temporal $b$-field tears that of one.

\begin{proposition}[Mass transfer between constituents]
\label{e6:prop:transfer}
Let $H_g=m_1E_{(1,1)}+m_2E_{(2,1)}$, $M=(m_1^2+m_2^2)^{1/2}$, $u_1=e_1\otimes\id$ and $u_2=\id\otimes e_1$.
\begin{enumerate}[label=(\roman*),leftmargin=2.2em]
\item $\Ad(e^{\theta C_+})E_{(1,1)}=\cos\theta\,E_{(1,1)}+\sin\theta\,E_{(2,1)}$ and $\Ad(e^{\theta C_+})E_{(2,1)}=\cos\theta\,E_{(2,1)}-\sin\theta\,E_{(1,1)}$. Hence
$\Ad(e^{\theta C_+})H_g=m_1(\theta)E_{(1,1)}+m_2(\theta)E_{(2,1)}$, where $(m_1(\theta),m_2(\theta))$ is $(m_1,m_2)$ rotated through $\theta$. The spectrum stays $\{\pm M\}$, and
$(M,0)$ goes to $(M\cos\theta,M\sin\theta)$. $C_-$ commutes with $H_g$.
\item The ungraded $H_u=m_1u_1+m_2u_2$ has spectrum $\{\pm m_1\pm m_2\}$, each value $16$ times, and this multiset determines $\{|m_1|,|m_2|\}$. Rotating
$(m_1,m_2)$ inside the ungraded family changes the spectrum, from $\{\pm M\}$ at $(M,0)$ to $\{\pm M(\cos\theta+\sin\theta),\pm M(\cos\theta-\sin\theta)\}$.
$(\cos\theta\,u_1+\sin\theta\,u_2)^2=\id+\sin2\theta\,u_1u_2$, so for $\sin2\theta\neq0$ no invertible matrix conjugates $u_1$ into $\cos\theta\,u_1+\sin\theta\,u_2$,
whereas $(\cos\theta\,E_{(1,1)}+\sin\theta\,E_{(2,1)})^2=\id$. The product $u_1u_2$ of the two ungraded clocks squares to $\id$ and commutes with both.
\item $\Ad(e^{\theta C_+})$ keeps the spectrum of $H_u$, but it maps $u_1$ to $\cos\theta\,u_1+\sin\theta\,\omega\otimes e_1$. This lies in
$\Cl(3,3)\otimes\id+\id\otimes\Cl(3,3)$ only for $\sin\theta=0$.
\end{enumerate}
\end{proposition}

\begin{proof}
(i) For $i=2,3$, $E_{(k,1)}$ anticommutes with both factors of $E_{(1,i)}E_{(2,i)}$, so these terms commute with $E_{(k,1)}$. The term
$X=E_{(1,1)}E_{(2,1)}$ anticommutes with $E_{(1,1)}$ and $E_{(2,1)}$, so $e^{-\theta X/2}E_{(1,1)}e^{\theta X/2}=e^{-\theta X}E_{(1,1)}=\cos\theta\,E_{(1,1)}+\sin\theta\,E_{(2,1)}$.
$C_-$ commutes with $E_{(k,1)}$ by the same count. (ii) $u_1$ and $u_2$ are commuting involutions whose products have trace zero, so each joint sign
occurs on $16$ dimensions. The largest and smallest absolute values in the multiset are $|m_1|+|m_2|$ and $||m_1|-|m_2||$. Conjugation preserves the
spectrum, and $\{\pm\cos\theta\pm\sin\theta\}\neq\{\pm1\}$ for $\sin2\theta\neq0$. (iii) $\omega\otimes e_1$ is a tensor of two non-identity blades, independent of every
$A\otimes\id$ and $\id\otimes B$.
\end{proof}

This is fixed (e6-6b, e6-6d). The checks are exact at $\cos\tfrac\theta2=\tfrac35$, where the rotated ungraded spectrum is
$\{\pm\tfrac{17}{25},\pm\tfrac{31}{25}\}M$. Part (i) is \cref{e1:prop:rest}(ii) realised by the temporal copy field. The single bivector $B=E_{(2,1)}E_{(1,1)}$
used there is one of the three terms of $2C_+$, and the other two leave the clocks alone.

\begin{remark}[Does a second particle acquire mass from the first?]
\label{e6:rem:transfer}
The algebra is fixed. The temporal copy rotation turns the clock $E_{(1,1)}$ of the first constituent into the clock $E_{(2,1)}$ of the second, and
from $(M,0)$ it produces $(M\cos\theta,M\sin\theta)$ at fixed graded mass: a massless second copy acquires mass from the first. Reading this as ``a
second particle acquires mass from the first through the $b$-field'' needs two further structures. One is a clock per copy, $E_{(k,1)}$ as the
$\gamma^0$ of copy $k$. The other is the graded reading, and there \cref{e1:prop:rest} applies: the copies have no joint mass shell, and the rest energy is
the single value $M$. In that reading the process is a rotation of the mass direction of one particle between two internal clock directions,
the masses adding in quadrature as in Proposition~8.11(iv) of \cite{RE}. It is not the creation of a second massive particle. In the ungraded
reading, where energies add and there are two particles (\cref{e1:prop:rest}(iii)--(iv)), the transfer is obstructed: no conjugation turns one
clock into the other, and the additive spectrum fixes the pair of masses (\cref{e6:prop:transfer}(ii)). The transfer is moreover not
binding-compatible, and half of it creates quanta in pairs across the constituents (\cref{e6:prop:copyrot}(iii)). So the mass transfer between
constituents is fixed algebra; the reading of it as the emergence of a second massive particle is interpretive in the one-particle reading and
obstructed in the two-particle reading.
\end{remark}

\subsection{Up the ladder: the mass forms of a rung}
\label{e6:sec:ladder}

\begin{theorem}[Mass forms of a rung]
\label{e6:thm:ladder}
At rung $n$ let $T_n$ be the product, in increasing order, of the $3n$ temporal generators $E_{(k,i)}$ ($i\le3$), $S_n$ that of the $3n$ spatial ones, and
$\Gamma_n$ the volume element of $\Cl(3n,3n)$.
\begin{enumerate}[label=(\roman*),leftmargin=2.2em]
\item For every $n$, $Y^{\!\top}T_n+T_nY=0=Y^{\!\top}S_n+S_nY$ for every bivector $Y$, $T_nS_n=\pm\Gamma_n$, and the forms on $S_n$ invariant under $\so{3n,3n}$
are exactly $\operatorname{span}\{T_n,S_n\}$.
\item $T_n^{\!\top}=(-1)^{3n(3n-1)/2}T_n$ and $S_n^{\!\top}=(-1)^{3n+3n(3n-1)/2}S_n$. For $n\equiv1,2,3,0\pmod4$ the pair $(T_n,S_n)$ is (antisymmetric,
symmetric), (antisymmetric, antisymmetric), (symmetric, antisymmetric), (symmetric, symmetric). Hence the Grassmann mass kernel invariant under
$\so{3n,3n}$ is unique iff $n$ is odd. It is $T_n$ for $n\equiv1$ and $S_n$ for $n\equiv3\pmod4$. For $n\equiv2$ the kernels form the plane $\operatorname{span}\{T_n,S_n\}$,
and for $n\equiv0$ there is none.
\item At $n=2$, $T_2=-\omega_S\otimes T$ and $S_2=-T\otimes\omega_S$. The four products $M_1\otimes M_2$, $M_k\in\{T,\omega_S\}$, span the forms invariant under
$\so{3,3}^{(1)}\oplus\so{3,3}^{(2)}$; $T\otimes T$ and $\omega_S\otimes\omega_S$ are symmetric, and $\so{6,6}$ keeps exactly the antisymmetric ones,
$T_2$ and $S_2$. The per-copy $T\otimes\id$ is not invariant even under the boosts of copy~$2$.
\item The forms invariant under the temporal rotations $\so{3n}_+$ are their commutant, of dimension $16$, $128$ and $1024$ at $n=1,2,3$, with
(symmetric, antisymmetric) dimensions $(6,10)$, $(64,64)$ and $(528,496)$. They contain $T_n$ and $S_n$.
\item The forms invariant under the diagonal $\so{3,3}$, the same bivector on every constituent, have dimension $2$, $14$ and $168$ at $n=1,2,3$, with
(symmetric, antisymmetric) dimensions $(1,1)$, $(10,4)$ and $(84,84)$.
\end{enumerate}
\end{theorem}

\begin{proof}
(i) At every rung the temporal generators are symmetric and the spatial ones antisymmetric, $E_A^{\!\top}=\eta_{AA}E_A$, since $\omega^{\!\top}=\omega$. By the
blade rule in $\Cl(3n,3n)$, $T_nE_AT_n^{-1}=(-1)^{3n-1}\eta_{AA}E_A=(-1)^{3n-1}E_A^{\!\top}$ for every $A$, so for $Y=E_AE_B$,
$T_nYT_n^{-1}=E_A^{\!\top}E_B^{\!\top}=(E_BE_A)^{\!\top}=-Y^{\!\top}$, which is the invariance. The same holds for $S_n$ with the sign $(-1)^{3n}$. Conversely, if $M$ is
invariant, then $Y^{\!\top}=-T_nYT_n^{-1}$ gives $Y(T_n^{-1}M)=(T_n^{-1}M)Y$ for every bivector. The commutant of the bivectors is spanned by the blades it
contains. A blade $e_I$ commutes with $E_AE_B$ iff $|I\cap\{A,B\}|$ is even, and this holds for all $A\neq B$ iff $I=\emptyset$ or $I$ is everything. So
$M\in T_n\operatorname{span}\{\id,\Gamma_n\}=\operatorname{span}\{T_n,S_n\}$. (ii) The transpose of a product of $m$ symmetric (respectively antisymmetric)
generators is the reversed product times $1$ (respectively $(-1)^m$), and reversal costs $(-1)^{m(m-1)/2}$. (iii) $E_{(2,1)}E_{(2,2)}E_{(2,3)}=\omega^3\otimes T=\omega\otimes T$
and $T\omega=-\omega_S$. Within each copy the invariant forms are $\operatorname{span}\{T,\omega_S\}$, and blade conjugation is diagonal. (iv) A compact $Y$
is antisymmetric, so invariance is commutation; the dimensions are exact orbit counts. (v) Let $c(\mu)$ be the weight multiplicities of a module for
the split Cartan subalgebra spanned by the pole operators $e_ie_{i+3}$, which are diagonal. The multiplicity of the trivial module is
$\sum_{w\in W(D_3)}\det(w)\,c(\rho-w\rho)$ \cite{Humphreys72}. It is evaluated on $\mathrm{Sym}^2S_n$ and $\Lambda^2S_n$ with the exact weights; at $n=2$ a
direct exact solution on the $216$ entries of opposite weight agrees.
\end{proof}

The theorem is fixed (e6-7a to e6-7e). At $n=1$, (i) restates Corollary~18.4 of \cite{ES17} and the pairing of Section~28 there. The pattern of
(ii) is Chevalley's: on the spinors of $\R^{m,m}$ the two invariant forms have symmetries $(-1)^{m(m-1)/2}$ and $(-1)^{m(m+1)/2}$
\cite{Chevalley54,LawsonMichelsohn89}, here with $m=3n$ (inherited). What the arena adds is that the two forms are the temporal and the spatial volume
of the rung, with the Klein factors of (iii). The uniqueness of the mass in $T=e_1e_2e_3$ is therefore a statement of the first rung. At the paper's
three-constituent rung the unique $\so{9,9}$-invariant Grassmann kernel is the spatial volume $S_3$, and $T_3$ is a symmetric (bosonic) form. Under
the diagonal action the mass is not unique from $n=2$ on, and under the temporal rotations alone not even at $n=1$. At the second rung the
antisymmetric invariant kernels pair the temporal volume of one constituent with the spatial volume of the other (iii): no invariant mass kernel
belongs to one constituent alone.

\subsection{The two assignments compared}
\label{e6:sec:compare}

\begin{remark}[The temporal $b$-field and the $b$-field of \cite{RE}]
\label{e6:rem:compare}
\begin{enumerate}[label=(\alph*),leftmargin=2.2em]
\item \emph{How they relate} (fixed). In Sections~5 and~9 of \cite{RE} the $b$-field is the $\eta$-antisymmetric part of the shift $S=\kappa G_n+b$, and
$\kappa G_n$ is its $\eta$-self-adjoint part, $(\kappa G)^2=-c^2/\ell^2$; $\kappa G$ maps $V^+$ to $V^-$ and back, and that is where $c$ enters. Section~9 of
\cite{RE} studies the binding-compatible fields, $[b,G_n]=0$. A $b$-field on the temporal triple is part of the connection, compact and
gauge-admissible (Corollary~2.4 of \cite{ES17}), but it is not binding-compatible: half of it lies in $\mathfrak p_G$
(\cref{e6:prop:binding}(i)). It belongs to the compact class of Section~5 of \cite{RE} with $\mathbf a=\mathbf c$, it transports $\kappa G$ to
$\kappa G_{e^{-\theta R}}$ without changing the rates, and it tears the binding (\cref{e6:prop:binding}). Between constituents the same holds for the
temporal copy rotation (\cref{e6:prop:copyrot}(iv)), which lies outside the Schr\"odinger class of Theorem~9.3 of \cite{RE}.
\item \emph{A non-abelian Yang--Mills $\su2$ with Cartan direction $J$} (fixed; the Cartan choice is the clock's). $\so3_+\cong\su2$ is compact. It acts on
$V$ as isospin one on $V^+$ and trivially on $V^-$, and on spinors as isospin one-half (\cref{e6:prop:dirac}(iii)). The algebra fixes no Cartan direction,
since all elements of $\so3_+$ are conjugate. The kinetic term fixes it to $J=J(e_1)$, the rotation about the clock, which it keeps. The two
off-diagonal directions are the clock rotations; they flip the charge, and their half turn acts as $\mathsf{CP}$.
\item \emph{A mass term for the off-diagonal components} (identity fixed; reading interpretive; mechanism open). The breaking of the kinetic forms
by $Y\in\so3_+$ is $Y^{\!\top}K^a+K^aY=2Y(n)\gamma^a$, with total weight $96|Y(n)|^2$ (\cref{e6:prop:dirac}(i)). This is a quadratic form on $\so3_+$ with
kernel $\R J(n)$, invariant under the stabiliser of the clock: an invariant statement of the breaking. It has the form of the adjoint-Higgs mass
term $g^2|[b,\langle\Phi\rangle]|^2$ of the Georgi--Glashow $SO(3)$ model \cite{GeorgiGlashowPRL72}, with the clock in the role of the Higgs vector
\cite{EnglertBrout64,Higgs64,Kibble67}: the two charge-flipping components would be massive and $J(n)$ massless. That needs a clock field $n(x)$ with a
covariant kinetic term, and the arena supplies none: a dynamics that selects the clock is open (Section~3 of \cite{TR}). Gauge invariance forbids a
bare mass term (inherited), and the Dirac mass gives none, since $\so3_+$ preserves $T$ (\cref{e6:prop:Tform}). In Remark~9.8(e) of \cite{RE} the rate matrix
played the role of the adjoint vacuum value; here the clock does.
\item \emph{What the assignment fixes, constructs, obstructs and leaves open.} Fixed: every finite-dimensional statement of
\cref{e6:sec:Tform,e6:sec:dirac,e6:sec:binding,e6:sec:polar,e6:sec:clock,e6:sec:copies,e6:sec:ladder}. Constructed: the clock, which fixes the Cartan
direction and the kinetic term, and the temporal $b$-field itself. Obstructed: its removal along a worldline, bounded binding under it, mass transfer
at fixed spectrum in the ungraded reading, and a unique Grassmann mass kernel at even rungs. Open: a dynamics for the clock, hence a mass for the
off-diagonal components.
\item \emph{What it implies for the emergence of a second particle} (interpretive, drawn from the fixed statements). The temporal $b$-field gives the
second rung a process that moves mass from one constituent to the other, but only as a rotation of one mass vector in the graded reading, in which
the rung carries one particle. In the reading that has two particles it is obstructed. The uniqueness of the mass in $T$, which the assignment invokes,
does not survive the ladder in a form that assigns a mass to each particle: the invariant kernels of the second rung pair the temporal volume of one
constituent with the spatial volume of the other.
\end{enumerate}
\end{remark}

\subsection{Answer}
\label{e6:sec:answer}

The assignment ``$c$ on $(-,-,-)$, $b$ on $(+,+,+)$'' has a precise and correct form. The temporal rotations cannot change the symmetric factor
$CC^{\!\top}$ of the conversion, which lives on the spatial triple and is $c^2$ in the isotropic case; they keep the mass form $T$ and the rest
flow, and they move the orthogonal factor of the conversion, the clock and the binding root. The price is the binding and the kinetic term: a
temporal $b$-field keeps the rates, splits the lines, cannot be gauged away and tears the binding, and the kinetic term breaks it to $\uu1_J$ in
proportion to the speed of the clock. The mass is unique in $T$ on the first rung only. Between two constituents the temporal $b$-field rotates
graded mass from one clock to the other at fixed total, a statement about one particle, not two.

\begin{table}[htbp]
\centering
\footnotesize
\begin{tabular}{@{}>{\raggedright\arraybackslash}p{8.0cm}>{\raggedright\arraybackslash}p{4.7cm}>{\raggedright\arraybackslash}p{2.1cm}@{}}
\toprule
statement & class & checks\\
\midrule
conventions; the triples, $T$, $\omega_S$, Fock modes & restated (ES17 eqs.~(1), (2), Thm~18.1) & e6-0\\
\cref{e6:prop:Tform}(i): $T$ and $\omega_S$ invariant forms under $\mathrm{Spin}(3,3)$ & restated (ES17 Cor.~18.4, Sec.~28; RE Prop.~3.5) & e6-1a\\
\cref{e6:prop:Tform}(ii), (iii): commutes with $\so3_+\oplus\so3_-$ only; commutant $M_4(\C)$, $32$ blades & fixed ($M_4(\C)$ inherited) & e6-1b\\
\cref{e6:prop:Tform}(iv): $\Ad_T=\eta$; boosts conjugate $T$ to $e^{2sY}T$ & restated (RE Prop.~2.11); fixed & e6-1c\\
\cref{e6:prop:dirac}(i): kinetic breaking $2Y(n)\gamma^a$, weight $96|Y(n)|^2$; $\so3_+\cap(\lint\oplus\R J)=\R J$ & fixed (covariance restated, RE Prop.~3.5(ii)) & e6-2a\\
\cref{e6:prop:dirac}(ii): clock rotation of $H(0)$; transport to the clock $n$ & fixed (transport restated, TR Prop.~3.4, Thm~3.5) & e6-2b\\
\cref{e6:prop:dirac}(iii), (iv): rest flow kept; isospin $\tfrac12$; $\mathsf{CP}$ half turn; octant duality & fixed; (iv) restated (TR Prop.~4.12) & e6-2c\\
\cref{e6:prop:twosu2}: $\so3_+$ and $\su2_{\rm PG}$ two diagonals of $\so4$, meeting in $\R J$ & fixed ((v) restated, ES17 Prop.~5.5) & e6-2d\\
\cref{e6:prop:binding}(i), (ii): half $\lG$, half $\mathfrak p_G$; $G\to G_{e^{-\theta R}}$; rates kept; not removable & fixed (rates restated, RE Thm~9.3(iii)) & e6-3a, e6-3b\\
\cref{e6:prop:binding}(iii), (iv): lines $\kappa$, $s\pm\omega/2$; Jordan blocks & fixed & e6-3c, e6-3d\\
\cref{e6:prop:binding}(v): paired half conserves $N$, counter half changes it by $\pm2$ & fixed & e6-3e\\
\cref{e6:prop:polar}: $C=cR$; twisted diagonal; orbit $O(3)$; $CC^{\!\top}$ invariant; $q_C$ & fixed & e6-4a to e6-4e\\
\cref{e6:rem:csym}: ``$c^2$ on $(-,-,-)$, $b$ on $(+,+,+)$'' & (a), (b) fixed; reading interpretive & ---\\
\cref{e6:prop:clock}: $m_4$ oscillates at fixed $\mu$; constant iff on the axis & fixed; the reading interpretive & e6-5a, e6-5b\\
\cref{e6:prop:copyrot}: $C_\pm$, $e^{2\pi C_\pm}=-\id$; pairs of quanta; tearing between frames & fixed ((i) restated, RE Prop.~9.6(i)) & e6-6a, e6-6c, e6-6e\\
\cref{e6:prop:transfer}: graded mass transfer at fixed $M$; ungraded obstruction & fixed & e6-6b, e6-6d\\
\cref{e6:rem:transfer}: ``a second particle acquires mass'' & interpretive (one-particle reading); obstructed (two-particle reading) & ---\\
\cref{e6:thm:ladder}(i), (ii): invariant forms $\operatorname{span}\{T_n,S_n\}$; period four & fixed (Chevalley's pattern inherited) & e6-7a, e6-7b\\
\cref{e6:thm:ladder}(iii)--(v): Klein factors; temporal and diagonal counts & fixed & e6-7a, e6-7c to e6-7e\\
\cref{e6:rem:compare}: relation to RE; Yang--Mills $\su2$; clock as adjoint Higgs & (a), (b) fixed; (c) identity fixed, reading interpretive, mechanism open; (e) interpretive & e6-2a\\
\bottomrule
\end{tabular}
\caption{Claims of \cref{e6:sec:main}, their classes and their checks, which are \texttt{PASS} lines of \texttt{verify/e6\_bfield\_triples.py}. ``Restated'':
fixed in \cite{ES17}, \cite{RE} or \cite{TR}, re-verified here.}
\label{e6:tab:claims}
\end{table}
